\section{Implementation in the BX3 Framework}\label{sec:bailout-implementation}

The preceding sections specify the Bailout Protocol as a formal state machine
$\mathcal{B} = (Q, q_0, \Sigma, \rightarrow, Q_F)$ and as a compulsory escalation
mechanism operating across a recursive agent tree.  This section shows how that
specification is realized within the BX3 three-layer architecture: the Purpose
Layer as the exclusive human accountability anchor, the Bounds Engine as the
non-discretionary trigger enforcer, the Fact Layer as the enforcement substrate
and Forensic Ledger authority, and the state machine itself embedded as a
deterministic Fact Layer extension that makes BX3's upstream accountability
guarantee structurally operative at runtime.

% --------------------------------------------------------------
\subsection{Purpose Layer: The Human Accountability Anchor}
\label{sec:bailout-purpose-layer}

The Purpose Layer carries intent and final authority.  Its role in the Bailout
Protocol is singular and architecturally irreplaceable: it is the exclusive
terminus for all exception propagation and the only layer capable of issuing a
binding resolution directive that closes a protocol run.

Each BX3 node $n$ is provisioned at creation with a reference to a named human
principal $h(n)$.  This reference is stored in the node's Purpose Layer worksheet
and is immutable for the node's operational lifetime.  No machine actor---including
the node's own Bounds Engine, any parent node, or any recursively spawned
subordinate---may modify $h(n)$ without explicit human authorization recorded as
a ledger entry.  This immutability is not a convention; it is structurally
enforced by the Fact Layer's pre-commit hook (IC-4): any proposed modification
to $h(n)$ requires a new ledger entry authorized by the current $h(n)$ itself,
creating a circular authorization constraint that renders the field effectively
immutable in practice and architecturally immutable by design.

The Purpose Layer's participation in the Bailout Protocol proceeds through four
functional stages:

\begin{itemize}\itemsep=0pt
\item[(\textit{i})] \textbf{Provisioning.}  At node initialization, the Purpose
Layer registers $h(n)$ with the Fact Layer's Pathway Registry and records the
anchor assignment in the authorization ledger.  This entry constitutes the human
principal's explicit acceptance of accountability for node $n$'s operational
scope.  The anchor is non-negotiable and non-migratory: for all nodes $n_i$ in
any spawn chain $C$ originating from $n$, $h(n_i) = h(n)$.

\item[(\textit{ii})] \textbf{Exception receipt.}  When an exception state
propagates to the Purpose Layer via the Escalate algorithm, the Purpose Layer
receives the full diagnostic context accumulated along the propagation chain: the
originating node, the trigger condition $\text{TRIGGER}(n, x)$, the operating
range definition $R(n)$, the time since first detection, and the resolution
history of any prior occurrences of this exception class.  No machine actor may
truncate, aggregate, or redact this context before it reaches the Purpose Layer.

\item[(\textit{iii})] \textbf{Directive issuance.}  The Purpose Layer evaluates
the exception and issues one of three binding directives: \texttt{ACK}
(acknowledge and resume normal operation), \texttt{RESOLVE} (execute a specified
binding resolution and resume), or \texttt{REJECT} (disallow the proposed action
and enter a quiescent error state).  No machine actor may issue, simulate, or
substitute any of these directives.  Upon receipt of the directive, the state
machine transitions to \texttt{HUMAN\_ACKNOWLEDGED} and then to \texttt{RESOLVED};
the directive is committed to the authorization ledger as a Purpose-Layer-
authorized entry before it takes effect.

\item[(\textit{iv})] \textbf{Consequential authority.}  The Purpose Layer is the
only layer that carries the authority to commit consequential decisions to the
authorization ledger.  This exclusivity is what makes it the only permissible
terminus for the accountability chain: any attempt by a machine actor to close a
protocol run without a ledger entry bearing a Purpose-Layer-authorized directive
fails the Fact Layer's pre-commit hook and is recorded as a protocol violation.
\end{itemize}

The human anchor $h(n)$ is not a design option.  It is a structural necessity.
The functional separation between Purpose, Bounds Engine, and Fact Layers
establishes that the Purpose Layer is the only layer carrying both intent and
final accountability.  The Bounds Engine detects exception conditions but carries
no judgment authority; the Fact Layer enforces transition relations and records
events but carries no intent authority.  When an exception exceeds the Bounds
Engine's provisioned scope, no machine actor exists that has the authority to
adjudicate it.  The decision must return to the Purpose Layer, which carries the
intent and accountability that any such adjudication requires.  The Bailout
Protocol is the mechanism by which this structural necessity is enforced at
runtime.

The non-substitutability property follows directly: no machine actor operating
under delegated authority possesses the principal's intent in sufficient
completeness to substitute for $h(n)$ in the accountability chain.  The
principal's intent is defined at provisioning and carried exclusively in the
Purpose Layer; any machine actor's authority is inherently a strict subset of
that intent, operating within $R(n)$, and therefore cannot serve as a terminus for
exceptions that exceed its own scope.

% --------------------------------------------------------------
\subsection{Bounds Engine: Triggering Bail-Out}
\label{sec:bailout-bounds-engine}

The Bounds Engine carries bounded reasoning and constrained execution within the
provisioned operating range $R(n)$.  Its participation in the Bailout Protocol
is confined to a single, non-discretionary function: detecting conditions that
fall outside $R(n)$ and triggering the state machine's compulsory transition to
\texttt{BOUNDED}.

The Bounds Engine maintains a read-only view of $R(n)$, provisioned at node
startup and immutably protected by the Fact Layer's pre-commit hook.  On every
consequential action proposed by the node, the Bounds Engine evaluates whether
the action falls within $R(n)$ by querying the authorization ledger via the Fact
Layer (IC-1).  If the action is within $R(n)$, it proceeds through the standard
execution path.  If the action---or an observed exception condition $x$---falls
outside $R(n)$, the Bounds Engine evaluates the full trigger condition TC:

\begin{equation}
\text{TRIGGER}(n, x) \equiv x \notin R(n) \;\land\; \neg\text{resolved}(n, x) \;\land\;
\text{confidence}(n, x) < \tau .
\tag{TC}
\end{equation}

Critically, the Bounds Engine does not exercise discretion over this
evaluation.  It performs a deterministic read against the operating range
definition and the confidence function, both provisioned at node initialization
and immutable at runtime.  There is no judgment call: the trigger condition
either fires or it does not, and the result directly drives the state machine
transition without an intervening machine-actor gate.  The evaluation of
$\text{TRIGGER}(n, x)$ is driven directly by the Bailout Monitor's pre-commit
hook, which intercepts the query and commits the result to the state machine
before any agent-level code can act on the input.

When TC fires, the state machine transitions from \texttt{IDLE} to
\texttt{BOUNDED} (T1), and the Bounds Engine begins a time-bounded resolution
attempt with timeout $T_{\text{bounds}}$.  During the \texttt{BOUNDED} state, all
resolution paths within current authority remain open; the Bounds Engine may
attempt multiple strategies in sequence or in parallel, each submission reviewed
by the Fact Layer against the authorization ledger.  If the Fact Layer approves
a resolution before $T_{\text{bounds}}$ expires, the state machine returns to
\texttt{IDLE} (T2).  If $T_{\text{bounds}}$ expires without an approved
resolution, the state machine transitions to \texttt{BAIL\_PENDING} (T3) and the
Bounds Engine emits the \texttt{bailout\_signal} to the Escalate algorithm.

The no-discretion property of the trigger is the architectural feature that
prevents machine actors from suppressing warranted escalations.  Because the
trigger evaluation is intercepted and committed by the Bailout Monitor before any
agent-level code observes its result, no Bounds Engine instance can choose to
suppress a trigger condition on the grounds that it believes the exception can be
handled autonomously.  The determination that an exception exceeds the node's
authority is not a judgment the Bounds Engine is empowered to make; it is a
structural signal generated by the comparison of the observed condition against
the provisioned operating range.

% --------------------------------------------------------------
\subsection{Fact Layer: Enforcement and Forensic Ledger}
\label{sec:bailout-fact-layer}

The Fact Layer carries deterministic enforcement and forensic auditability.  Within
the Bailout Protocol, it serves two distinct but cooperating functions:
enforcing the state machine's transition relation, and maintaining the append-only
Forensic Ledger as an immutable record of every protocol event.

\medskip\noindent
\textbf{State machine enforcement.}  The Fact Layer enforces the Bailout Protocol
through three cooperating mechanisms:

\begin{itemize}\itemsep=0pt
\item[(\textit{i})] \textbf{Pre-commit hook.}  Before any state transition is
committed, the Fact Layer's Bailout Monitor evaluates whether the proposed
transition satisfies the current state's enabled transition set.  If no
transition is defined for the $(q, \sigma)$ pair, the transition is rejected and a
protocol violation is recorded in the Forensic Ledger.  This hook runs
synchronously and atomically: no intermediate or partially-executing state is
observable externally.

\item[(\textit{ii})] \textbf{Authorization ledger gate.}  The Fact Layer is the
sole write authority for the authorization ledger.  Any resolution that purports
to close a protocol run without a corresponding Purpose-Layer-authorized
directive fails the gate and is rejected.  This prevents the Bounds Engine from
issuing a retroactive resolution after timeout to suppress a warranted escalation.

\item[(\textit{iii})] \textbf{Pathway Registry.}  The Fact Layer maintains the
Pathway Health Registry, which tracks all communication pathways to the human
anchor $h(n)$.  When a pathway is marked degraded or unavailable, the registry
enables the Escalate algorithm to identify and route through alternative
functional pathways.  The registry is append-only for historical health data;
current-pathway state is a mutable field updated by the Bailout Monitor.
\end{itemize}

\medskip\noindent
\textbf{Forensic Ledger recording.}  The Forensic Ledger is the immutable
evidentiary record produced by the Fact Layer's enforcement of the Bailout
Protocol.  Every protocol event---trigger evaluations, state transitions,
pathway selections, acknowledgments, directives, and timeout conditions---is
committed to the ledger as an append-only entry.  Each entry is timestamped,
attributed to the originating node, and linked to its predecessor by a SHA-256
hash chain that makes retrospective tampering computationally detectable.

The Forensic Ledger records each protocol event across three discrete
accountability planes:

\begin{enumerate}\itemsep=0pt
\item \textbf{Purpose Plane.}  Records the authorization context at the time of
bailout: the originating node, its human anchor $h(n)$, the active mandate, and
the constraint boundary or trigger condition that was exceeded.  This plane
establishes the provenance of the exception in the delegation chain.

\item \textbf{Bounds Plane.}  Records the diagnostic context accumulated during
propagation: the trigger condition, each proposed resolution attempt, the
Bounds Engine's confidence level at time of firing, and the complete propagation
chain showing every node the exception traversed before reaching the Purpose
Layer.

\item \textbf{Fact Plane.}  Records the authoritative state transitions and
enforcement events: the Fact Layer's pre-commit hook decisions, the state machine
transitions committed, timeout expirations, and the human principal's directive
with its authorization ledger reference.  This plane provides the auditable
evidence that the protocol operated correctly at every step.
\end{enumerate}

The SHA-256 cryptographic chaining ensures that no entry can be modified,
deleted, or inserted without detection.  An auditor need only verify the chain
integrity by recomputing each hash and comparing it to the stored successor hash
to confirm that no tampering has occurred---without access to the ledger's
contents.  The complete Forensic Ledger record for a Bailout Protocol escalation
constitutes the definitive accountability artifact: it captures the full
provenance of the exception from the triggering input through every intermediate
node's resolution attempt to the human principal's binding directive.

% --------------------------------------------------------------
\subsection{Relationship to the Upstream Accountability Guarantee}
\label{sec:bailout-accountability}

The upstream accountability guarantee---that BX3 systems fail upward into human
consciousness, never downward into autonomous chaos---is realized through the
coordinated operation of all five BX3 pillars.  The Bailout Protocol is its
runtime expression: the mechanism that converts the static structural commitment
into a deterministic, auditable, architecturally enforced propagation event.

Loop Isolation, Recursive Spawning, Spatial Firewall, and Sandbox Gate address
exception conditions anticipated at design time and resolvable within the
machine-only execution graph.  The Bailout Protocol addresses the complementary
remainder---conditions not anticipated, that exceed the operating range of any
machine actor, and that cannot be resolved within the machine-only execution
graph.  The completeness of this coverage is what makes the upstream
accountability guarantee unconditional: it is not that \emph{most} exceptions
reach $h(n)$, but that \emph{every} exception reaches $h(n)$, because the trigger
condition fires on any condition outside $R(n)$ and the escalation path has no
machine-actor terminus by architectural constraint.

The relationship between the protocol and the guarantee is one of mutual
structural definition.  The upstream accountability guarantee requires the
Bailout Protocol to be operative: without it, the guarantee would be
unenforceable for the exception conditions that the four preceding pillars do not
cover.  The guarantee is a claim about the behavior of the system; the protocol
is the mechanism that makes the claim true.  Conversely, the Bailout Protocol
requires the structural properties that only the BX3 three-layer architecture
provides: the Purpose Layer must be non-delegable as an exception terminus, the
$h(n)$ anchor must be immutable at runtime, and the Fact Layer must have
exclusive write authority over the Forensic Ledger.  These are not optional
features that improve the protocol's performance---they are structural
requirements without which the protocol's guarantees cannot be realized.

The Forensic Ledger entries generated by the protocol provide the retrospective
proof of this guarantee.  Any external auditor---regulator, customer, investor,
or court---can reconstruct the complete exception history of any node from the
cryptographically chained Forensic Ledger entries and verify that every unresolved
exception reached $h(n)$, that every human directive was recorded, and that no
machine actor suppressed, redirected, or substituted for the human principal.

% --------------------------------------------------------------
\subsection{State Machine as Fact Layer Extension}
\label{sec:bailout-state-machine}

The Bailout Protocol state machine $\mathcal{B}$ is embedded in the Fact Layer
of every BX3 node as a Fact Layer extension---a deterministic enforcement
automaton whose transition relation is enforced by the Fact Layer's pre-commit
hook, whose state is stored in the Fact Layer's worksheet, and whose state
transitions generate entries in the Forensic Ledger as co-commitments.

This embedding is not incidental.  The state machine's enforcement requirements
are precisely matched to the Fact Layer's architectural properties:

\begin{itemize}\itemsep=0pt
\item[(\textit{i})] \textbf{Deterministic transition enforcement.}  For every
$(q, \sigma) \in Q \times \Sigma$, at most one transition is defined.  This
determinism is delivered by the Bailout Monitor's priority function $\pi$, which
totally orders simultaneous trigger conditions before committing any state change.
The Fact Layer provides the scheduling guarantee and atomic pre-commit hook
required to make this ordering consistent across all nodes.

\item[(\textit{ii})] \textbf{Atomic state transitions.}  No intermediate or
partially-executing state is observable externally.  The Fact Layer's pre-commit
hook evaluates and commits the transition in a single synchronized operation,
preventing any interleaving machine actor from observing or modifying the state
mid-transition.

\item[(\textit{iii})] \textbf{State observability.}  The state machine's current
state is observable by the parent node via the Fact Layer's state observation
interface, enabling parent-level monitoring of child escalation status
independently of the child node's internal representation.

\item[(\textit{iv})] \textbf{Ledger co-commitment.}  Every state transition
generates a mandatory Forensic Ledger entry as a co-commitment of the transition
itself.  No state change in $\mathcal{B}$ is committed without a corresponding
ledger entry recording the transition event, the triggering condition, and the
diagnostic context.
\end{itemize}

The state machine is therefore not merely \emph{located} in the Fact Layer---it
is defined by the Fact Layer's enforcement capabilities.  The six-state machine,
the seven-transition relation (T1--T7), the timeout parameters, and the
escalation algorithm are all realized as Fact Layer primitives.  The Fact Layer
does not merely host the state machine; it \emph{is} the state machine's
enforcement substrate.

The state invariants maintained by the Fact Layer's enforcement architecture are:

\begin{align}
\text{INV}_1 &: \forall n \in \text{nodes},\;
\text{state}(n) \in Q \setminus \{\text{ESCALATING}\}
\iff \neg\text{active\_bailout}(n). \tag{BP-INV-1}
\end{align}

INV-1 is enforced by the Fact Layer's pre-commit hook: a node may enter a
non-escalating state only when the Bailout Monitor confirms that no active
propagation is in flight.  The state and the active-bailout flag are updated
atomically, preserving the equivalence.

\begin{align}
\text{INV}_2 &: \text{state}(n) = \text{BAIL\_PENDING}
\implies \text{active\_bailout}(n) \land
\text{transition}(\text{BAIL\_PENDING}, e_{\text{timeout}}) = \text{ESCALATING}.
\tag{BP-INV-2}
\end{align}

INV-2 encodes the no-machine-actor-discretion property as a state machine
invariant.  The transition from \texttt{BAIL\_PENDING} to \texttt{ESCALATING} is
forced by the state machine's deterministic transition relation and enforced by
the Bailout Monitor without machine-actor intervention.  The Fact Layer's
pre-commit hook ensures this transition cannot be intercepted or redirected.

Together, the state invariants and the Fact Layer's enforcement architecture
ensure that the Bailout Protocol state machine operates as a deterministic,
atomic, auditable extension of the Fact Layer---and that every transition it
executes is immutably recorded in the Forensic Ledger as evidence that the
upstream accountability guarantee has been upheld.
